\section*{Resumen}

Este trabajo formula y evalúa un modelo dinámico de caja gris para la
generación de dióxido de carbono durante la bioconversión de residuos
orgánicos con \textit{Hermetia illucens}. El punto de partida es el
núcleo larvario de Eriksen, expresado en equivalentes de carbono, que
conserva los balances de biomasa estructural, lípidos, mantenimiento y
costos respiratorios de crecimiento. Sobre ese núcleo se incorporan dos
componentes: una fuente microbiana asociada con la degradación del
alimento y un acople entre la disponibilidad de alimento y la
asimilación larvaria. El dióxido de carbono medido en cámara se
representa como la suma de las contribuciones larvaria y microbiana.

En ocho cierres de gas de los experimentos propios, el núcleo original
sin reajuste produjo un RMSE conjunto de $183{,}61$ ppm/min. El modelo
completo recalibrado redujo el error a $85{,}16$ ppm/min, equivalente a
una disminución de $53{,}6$\,\%. La reducción no debe atribuirse
exclusivamente al lecho: el ajuste conjunto incluye la recalibración de
la asimilación larvaria, el acople con el alimento y la nueva fuente
microbiana.

La transferencia a Laganaro, Bekker, Ermolaev, Schøn y Rossi muestra
capacidades y límites diferentes en cada estudio. El modelo propuesto
mejora algunos errores agregados, pero no reproduce universalmente el
efecto de dietas de baja calidad, la distribución temporal de los picos,
la partición microbiana ni el colapso a temperaturas extremas. El
resultado principal es, por tanto, una mejora local en los datos propios
y una representación explícita de las fuentes de \co{}, no una
superioridad general frente al modelo de Eriksen.
